
\section{$n$ 个线丛与 $n$ 维簇的\trIntersection}

本章的中心工具是下面这个。

\textbf{20.1.1. 定义：\trIntersectionProduct}（intersection product）\textbf{，或\trIntersectionNumber}（intersection number）$(\mathscr{L}_1\cdots\mathscr{L}_n)$、$(\mathscr{L}_1\cdots\mathscr{L}_n\cdot\mathscr{F})$、$(\mathscr{L}_1\cdots\mathscr{L}_n\cdot Y)$。\quad 设 $\mathscr{F}$ 是 $X$ 上具有固有支撑的一个\trCoherentSheaf（若 $X$ 固有则自动如此），其\emph{维数至多为 $n$}，而 $\mathscr{L}_1,\dots,\mathscr{L}_n$ 是 $X$ 上的\trInvertibleSheaf。令 $(\mathscr{L}_1\cdot\mathscr{L}_2\cdots\mathscr{L}_n\cdot\mathscr{F})$ 为对 $\{1,\dots,n\}$ 的 $2^n$ 个子集所取的带号和：
\begin{equation*}\tag{20.1.1.1}\label{eq:intersection-product-def}
\sum_{\{i_1,\dots,i_m\}\subset\{1,\dots,n\}}(-1)^m\chi(X,\mathscr{L}_{i_1}^{\vee}\otimes\cdots\otimes\mathscr{L}_{i_m}^{\vee}\otimes\mathscr{F}).
\end{equation*}
我们称其为 $\mathscr{L}_1,\dots,\mathscr{L}_n$ 与 $\mathscr{F}$ 的\textbf{\trIntersection}（intersection）。（\emph{切勿遗忘}：凡写下 $(\mathscr{L}_1\cdots\mathscr{L}_n\cdot\mathscr{F})$，我们就已默认 $\dim\operatorname{Supp}\mathscr{F}\le n$。）

今后我们将发现，最有用的情形是 $\mathscr{F}$ 为某个\trClosedSubscheme{} $Y$（维数至多 $n$）的\trStructureSheaf{}的情形。此时可把它写作 $(\mathscr{L}_1\cdot\mathscr{L}_2\cdots\mathscr{L}_n\cdot Y)$，并称之为 $\mathscr{L}_1,\dots,\mathscr{L}_n$ 与 $Y$ 的\textbf{\trIntersection}。

若 $Y=X$，则可写作 $(\mathscr{L}_1\cdot\mathscr{L}_2\cdots\mathscr{L}_n)$，并称之为 $\mathscr{L}_1,\dots,\mathscr{L}_n$ 的\textbf{\trIntersection}（「$X$」隐而不书）。

若诸 $\mathscr{L}_i$ 全体相同，设为 $\mathscr{L}$，则人们常写作 $(\mathscr{L}^n\cdot\mathscr{F})$、$(\mathscr{L}^n\cdot Y)$ 或 $(\mathscr{L}^n)$。（当心这个令人混淆的记号：$\mathscr{L}^n$ 并不表示 $\mathscr{L}^{\otimes n}$。）

在某些情形下，按惯例可省去圆括号。

关于\trIntersectionProduct，本章将证明许多事情。有一件事实要留待我们研究\trFlatness{}之时（\mbox{练习~24.7.D}）：它是「形变不变的」——也就是说，它在「好」的族中恒常。

\textbf{20.1.A.（练习：现实检验）}\quad 证明：若 $\mathscr{L}_1\cong\OO_X$，则 $(\mathscr{L}_1\cdot\mathscr{L}_2\cdots\mathscr{L}_n\cdot\mathscr{F})=0$。

\textbf{20.1.B.（练习）}\quad 证明：\trIntersectionProduct \eqref{eq:intersection-product-def} 在 $k$ 的\trFieldExtension{}之下保持不变。

下面的练习表明，\trIntersectionProduct{}或许饶有趣味，因为它在两个有用的概念之间「内插」：其一是一条\trCurve{}上\trLineBundle{}的\trDegree，其二是\trBezout{}定理。

\textbf{20.1.C.（练习）}\quad（a）若 $X$ 是一条\trProjective{}\trCurve，而 $\mathscr{L}$ 是 $X$ 上的一个\trInvertibleSheaf，证明 $(\mathscr{L}\cdot X)=\deg_X\mathscr{L}$。

（b）设 $k$ 是无限域，$X=\mathbb{P}^N$，而 $Y$ 是 $X$ 的一个 $n$ 维\trSubvariety。若 $H_1,\dots,H_n$ 是泛取的\trHypersurface，\trDegree{}分别为 $d_1,\dots,d_n$（于是由\mbox{练习~12.3.D(d)}，$\dim(H_1\cap\cdots\cap H_n\cap Y)=0$），则由某个适当版本的\trBezout{}定理（不妨参看\mbox{练习~13.4.B} 或\mbox{练习~18.6.J}），
\[
\deg(H_1\cap\cdots\cap H_n\cap Y)=d_1\cdots d_n\deg(Y).
\]
证明
\[
(\OO_X(H_1)\cdots\OO_X(H_n)\cdot Y)=d_1\cdots d_n\deg(Y).
\]

现在我们来描述\trIntersectionProduct{}的若干性质。在证明\mbox{练习~20.1.C(b)} 的过程中，你实际上可能已顺带解出了下面的练习。

\textbf{20.1.D.（练习）}\quad 设 $D$ 是 $X$ 上的一个有效\trCartierDivisor，它在 $Y$ 上的限制是一个有效\trCartierDivisor{} $D|_Y=D\cap Y$（也就是说，它在 $Y$ 上局部仍不是\trZerodivisor）。证明
\[
(\mathscr{L}_1\cdots\mathscr{L}_{n-1}\cdot\OO(D)\cdot Y)=(\mathscr{L}_1\cdots\mathscr{L}_{n-1}\cdot D|_Y).
\]
更一般地，若 $D$ 是 $X$ 上的一个有效\trCartierDivisor，它不含 $\mathscr{F}$ 的任何\trAssociatedPoint，证明
\[
(\mathscr{L}_1\cdots\mathscr{L}_{n-1}\cdot\OO(D)\cdot\mathscr{F})=(\mathscr{L}_1\cdots\mathscr{L}_{n-1}\cdot\mathscr{F}|_D).
\]
（一个类似的想法曾在「\trHilbert{}多项式确实总是多项式」的证明中出现过；见 \eqref{eq:twisting-ses} 前后的讨论。）

\textbf{20.1.2.}\ \emph{\trhDefinition{}。}\quad 缘于此，当 $D$ 是一个有效\trCartierDivisor{}时，在\trIntersectionProduct{}的记号中，我们常写 $D$ 而不写 $\OO(D)$。我们把\trIntersection{}的双方随意地想成有效\trCartierDivisor，而不想成\trLineBundle。例如，我们将在 \S20.2 中讨论\trIntersectionTheory{}在曲面上的特殊情形；届时当我们\trIntersection{}两条\trCurve{} $C$ 与 $D$ 时，将把该\trIntersection{}写作 $(C\cdot D)$，甚至径直写作 $C\cdot D$。

\textbf{20.1.3.\ \trhProposition。}\ ——\ \emph{设 $X$ 是\trProjective{}的。对固定的 $\mathscr{F}$，\trIntersectionProduct{} $(\mathscr{L}_1\cdots\mathscr{L}_n\cdot\mathscr{F})$ 是诸 $\mathscr{L}_1,\dots,\mathscr{L}_n$ 的对称多重线性函数（symmetric multilinear function）。}

若把「射影」换成「固有」，\mbox{命题~20.1.3} 其实依然为真；参看 [Kl1, Prop.~2] 或 [Ko1, Prop.~VI.2.7]。与大多数向固有情形的推广不同，这并不只是\trChow{}引理的一次应用；它走的是另一条路子，用到了一个名为 \emph{d\'evissage}（「拆解」）的漂亮技巧。

\textbf{20.1.4. 重要注记：\trLineBundle{}的加性记号。}\quad 有一套标准、有用、却令人混淆的约定，它由\trIntersectionProduct{}的多重线性提示出来：我们有时把\trInvertibleSheaf{}的张量积\emph{加性地}书写。这在处理多重线性时颇为方便。（有些人为了免于混淆而改用除子而非线丛，因为当我们「乘」相应的线丛时，除子是在相加的。这在心理上有助益，却也可能添更多的乱：你还得分心去操心线丛究竟是否、为何、如何以及何时对应于除子\footnote{原书此处作「worry about the whether」，衍一 the，径删。——译者注}。）例如，在\mbox{练习~20.2.B--20.2.G}、\S20.2.11 与 \S20.4 中，我们就将使用这套约定。

\textbf{20.1.5.}\ \emph{\mbox{命题~20.1.3} 的证明。}\quad 对称性是显然的。由\mbox{练习~20.1.B}，我们可以假定 $k$ 是无限的（例如代数闭）。现在我们对 $n$ 作归纳来证明该结果。

\textbf{20.1.E.（练习：基础情形）}\quad 证明 $n=1$ 的情形。提示：\mbox{练习~18.4.M}。

现在我们假设该结果对更小的 $n$ 值成立。我们要用一个技巧。我们希望证明（对任意的 $\mathscr{L}_1,\mathscr{L}'_1,\mathscr{L}_2,\dots,\mathscr{L}_n$）：
\begin{equation*}\tag{20.1.5.1}\label{eq:additivity-expression}
(\mathscr{L}_1\cdot\mathscr{L}_2\cdots\mathscr{L}_n\cdot\mathscr{F})+(\mathscr{L}'_1\cdot\mathscr{L}_2\cdots\mathscr{L}_n\cdot\mathscr{F})-((\mathscr{L}_1\otimes\mathscr{L}'_1)\cdot\mathscr{L}_2\cdots\mathscr{L}_n\cdot\mathscr{F})
\end{equation*}
为 $0$。

\textbf{20.1.F.（练习）}\quad 把 \eqref{eq:additivity-expression} 改写为
\begin{equation*}\tag{20.1.5.2}\label{eq:rewrite-expression}
(\mathscr{L}_1\cdot\mathscr{L}'_1\cdot\mathscr{L}_2\cdots\mathscr{L}_n\cdot\mathscr{F}).
\end{equation*}
（此时乘积中出现了 $n+1$ 个线丛，但由于 $\dim\operatorname{Supp}\mathscr{F}\le n<n+1$，这并不与\trIntersectionProduct{}的定义相抵触。）

\textbf{20.1.G.（练习）}\quad 用归纳假设证明：若 $\mathscr{L}_n\cong\OO(D)$，其中 $D$ 是一个不含 $\mathscr{F}$ 的诸\trAssociatedPoint{}的有效\trCartierDivisor，则 \eqref{eq:additivity-expression} 为 $0$。

特别地，若 $\mathscr{L}_n$ 是\trVeryAmple{}的，则 \eqref{eq:additivity-expression} 为 $0$，因为\mbox{练习~18.6.A} 表明：$\mathscr{L}_n$ 存在一个在 $\mathscr{F}$ 的诸\trAssociatedPoint{}处不取零的\trSection。（当心：\mbox{练习~18.6.A} 要求域是无限的。\footnote{原书此处作「requres」，当为「requires」之误，径改。——译者注}若 $k$ 是有限的，你会怎样修改它使之奏效？提示：证明\mbox{练习~18.6.A} 的一个变体：对某个 $N\gg 0$，$\mathscr{L}_n^{\otimes N}$ 存在一个在 $\mathscr{F}$ 的诸\trAssociatedPoint{}处不取零的\trSection。对此的提示：比较\mbox{练习~12.3.D(c)} 与\mbox{练习~12.3.D(d)}。）

而凭它在 \eqref{eq:rewrite-expression} 中的化身所具有的对称性，表达式 \eqref{eq:additivity-expression} 在 $\mathscr{L}_1$ 是\trVeryAmple{}的时也消失。在证明的最后一个方程展示中，我们将反复使用这一观察。

由\mbox{练习~16.2.F}，$X$ 上的\emph{任何}\trInvertibleSheaf{}都可以写成「两个极丰富线丛之差」，故写 $\mathscr{L}_1=\mathscr{A}\otimes\mathscr{B}^{\vee}$、$\mathscr{L}'_1=\mathscr{A}'\otimes(\mathscr{B}')^{\vee}$，其中 $\mathscr{A}$、$\mathscr{B}$、$\mathscr{A}'$ 与 $\mathscr{B}'$ 都是\trVeryAmple{}的。那么
\begin{align*}
((\mathscr{L}_1\otimes\mathscr{L}'_1)\cdot\mathscr{L}_2\cdots\mathscr{L}_n\cdot\mathscr{F})\\
&= ((\mathscr{A}\otimes\mathscr{A}'\otimes(\mathscr{B}\otimes\mathscr{B}')^{\vee})\cdot\mathscr{L}_2\cdots\mathscr{L}_n\cdot\mathscr{F})\\
&=((\mathscr{A}\otimes\mathscr{A}')\cdot\mathscr{L}_2\cdots\mathscr{L}_n\cdot\mathscr{F})\\
&\hphantom{{}={}}-((\mathscr{B}\otimes\mathscr{B}')\cdot\mathscr{L}_2\cdots\mathscr{L}_n\cdot\mathscr{F})\\
&=(\mathscr{A}\cdot\mathscr{L}_2\cdots\mathscr{L}_n\cdot\mathscr{F})+(\mathscr{A}'\cdot\mathscr{L}_2\cdots\mathscr{L}_n\cdot\mathscr{F})\\
&\hphantom{{}={}}-(\mathscr{B}\cdot\mathscr{L}_2\cdots\mathscr{L}_n\cdot\mathscr{F})-(\mathscr{B}'\cdot\mathscr{L}_2\cdots\mathscr{L}_n\cdot\mathscr{F})\\
&=((\mathscr{A}\otimes\mathscr{B}^{\vee})\cdot\mathscr{L}_2\cdots\mathscr{L}_n\cdot\mathscr{F})\\
&\hphantom{{}={}}+((\mathscr{A}'\otimes(\mathscr{B}')^{\vee})\cdot\mathscr{L}_2\cdots\mathscr{L}_n\cdot\mathscr{F})\\
&=(\mathscr{L}_1\cdot\mathscr{L}_2\cdots\mathscr{L}_n\cdot\mathscr{F})+(\mathscr{L}'_1\cdot\mathscr{L}_2\cdots\mathscr{L}_n\cdot\mathscr{F}),
\end{align*}
正如所愿。\hfill$\square$

这份证明还附带给我们一份额外的收获。

\textbf{20.1.H.（练习）}\quad 设 $X$ 是\trProjective{}的。证明：若 $\dim\operatorname{Supp}\mathscr{F}<n+1$，而 $\mathscr{L}_1,\mathscr{L}'_1,\mathscr{L}_2,\dots,\mathscr{L}_n$ 是 $X$ 上的\trInvertibleSheaf，则 \eqref{eq:rewrite-expression} 为零。\emph{换言之，$n+1$ 个\trInvertibleSheaf{}与一个\trCoherentSheaf{} $\mathscr{F}$ 的\trIntersectionProduct，当 $\dim\operatorname{Supp}\mathscr{F}<n+1$ 时为零\footnote{原书此处作「vanishes if the dim Supp F」，衍一 the，径删。——译者注}。}（事实上，「射影」换成「固有」后结果依然成立，因为它所依赖的那些结果在这种更一般的情形下成立。）

\textbf{20.1.6.\ \trhProposition。}\ ——\ \emph{设 $X$ 是\trProjective{}的。\trIntersectionProduct{}只依赖于诸 $\mathscr{L}_i$ 的\trNumericallyEquivalent{}类。}

（\trNumericallyEquivalent{}定义于 \S18.4.12——两个\trLineBundle{}是\trNumericallyEquivalent{}的，若它们的\trDegree{}在任何\trCurve{}上都相同。\footnote{原书此处作「two line bundles are linearly equivalent」，当为「numerically equivalent」（数值等价）之误（\S18.4.12 所定义者即数值等价），径改。——译者注}）同样，「射影」换成「固有」后结果依然为真。但无论是在\mbox{命题~20.1.3} 的证明中，还是在\mbox{命题~20.1.6} 自身的证明中，我们都用到了「每个\trLineBundle{}都是两个\trVeryAmple{}\trLineBundle{}之差」这一事实。命题在固有情形下的证明，见 [FGIKNV, Prop.~B.20]。

\emph{证明。}\quad 设 $\mathscr{L}_1$ 与 $\mathscr{L}'_1$ \trNumericallyEquivalent，而 $\mathscr{L}_2,\dots,\mathscr{L}_n$ 与 $\mathscr{F}$ 任意。我们希望证明 $(\mathscr{L}_1\cdot\mathscr{L}_2\cdots\mathscr{L}_n\cdot\mathscr{F})=(\mathscr{L}'_1\cdot\mathscr{L}_2\cdots\mathscr{L}_n\cdot\mathscr{F})$。我们对 $n$ 作归纳来进行。$n=1$ 的情形由\mbox{练习~18.4.L} 得出。设 $n>1$，并设结果对「更小的 $n$」成立。由\mbox{练习~16.2.K}，我们可以写 $\mathscr{L}_n\cong\OO(D_1-D_2)$，其中 $D_1$ 与 $D_2$ 是不含 $\mathscr{F}$ 的诸\trAssociatedPoint{}的有效\trCartierDivisor。由\trIntersectionProduct{}的多重线性，只须就 $\mathscr{L}_n=\OO(D)$（其中 $D$ 是一个不含 $\mathscr{F}$ 的诸\trAssociatedPoint{}的有效\trCartierDivisor）的情形证明该结果。那么
\begin{align*}
(\mathscr{L}_1\cdot\mathscr{L}_2\cdots\mathscr{L}_n\cdot\mathscr{F}) &= (\mathscr{L}_1\cdot\mathscr{L}_2\cdots\mathscr{L}_{n-1}\cdot\mathscr{F}|_D) && \text{（\mbox{练习~20.1.D}）}\\
&= (\mathscr{L}'_1\cdot\mathscr{L}_2\cdots\mathscr{L}_{n-1}\cdot\mathscr{F}|_D) && \text{（归纳假设）}\\
&= (\mathscr{L}'_1\cdot\mathscr{L}_2\cdots\mathscr{L}_n\cdot\mathscr{F}) && \text{（\mbox{练习~20.1.D}）}.
\end{align*}
\hfill$\square$

\textbf{20.1.7.}\ \emph{渐近\trRiemannRoch。}\quad 若 $Y$ 是一条\trProjective{}（＝\trProper；见\mbox{练习~19.2.I}\footnote{原书此处作「Remark 19.2.I」，当为「Exercise 19.2.I」（练习~19.2.I）之误，径改。——译者注}）的\trCurve，则 $\chi(Y,\mathscr{L}^{\otimes m})=m\deg_Y\mathscr{L}+\chi(Y,\OO_Y)$ 是关于 $m$ 的一次多项式（见\mbox{练习~18.4.L} 与\mbox{练习~18.4.M}），其首项（leading term）是一个\trIntersectionProduct。这可以推广如下。

\textbf{20.1.I.（练习：渐近\trRiemannRoch）}\quad 设 $X$ 是\trProjective{}的。（照例，把「射影」换成「固有」后结果仍将成立。）设 $\mathscr{F}$ 是 $X$ 上满足 $\dim\operatorname{Supp}\mathscr{F}\le n$ 的一个\trCoherentSheaf，而 $\mathscr{L}$ 是 $X$ 上的一个\trLineBundle。证明：$q(m):=\chi(X,\mathscr{L}^{\otimes m}\otimes\mathscr{F})$ 是关于 $m$ 的、次数至多为 $n$ 的多项式。证明 $q(m)$ 中 $m^n$ 的系数（「首项」）是 $(\mathscr{L}^n\cdot\mathscr{F})/n!$。提示：\mbox{练习~20.1.H} 蕴涵
\[
(\underbrace{\mathscr{L}\cdot\mathscr{L}\cdots\mathscr{L}}_{n+1\text{ 次}}\cdot(\mathscr{L}^{\otimes i}\otimes\mathscr{F}))=0,
\]
而这进而（利用 \eqref{eq:intersection-product-def}）蕴涵
\[
\sum_{j=0}^{n+1}(-1)^j\binom{n+1}{j}q(m+j)=0
\]
对一切 $m$ 成立。你的论证可能与关于\trHilbert{}多项式之多项式性的证明（\mbox{定理~18.6.1}）相像，所以你或许能在那里找到进一步的提示。\mbox{练习~18.6.B} 尤其可能帮上忙。

我们知道这个多项式的全部系数——当 $X$ 是一条\trCurve{}时，由\trRiemannRoch{}定理（见 \eqref{eq:degree-coherent-sheaf}），或者基本上由定义便可知道。当 $X$ 是正则射影曲面而 $\mathscr{F}$ 是一个\trInvertibleSheaf{}时，等我们证明了曲面的\trRiemannRoch{}定理（\mbox{练习~20.2.B(b)}），我们也将知道／诠释全部系数。要理解一般的情形，我们需要\trChern{}类理论。其成果便是 \trHirzebruch--\trRiemannRoch{}定理，它还能被进一步推广为鼎鼎大名的 \trGrothendieck--\trRiemannRoch{}定理（见 [F2, \S15.2]）。

\textbf{20.1.J.（练习：\trProjection{}公式）}\quad 设 $\pi\colon X_1\to X_2$ 是两个同为 $n$ 维的\trIntegral{}射影\trVariety{}之间的一个（\trProjective{}）\trMorphism，而 $\mathscr{L}_1,\dots,\mathscr{L}_n$ 是 $X_2$ 上的\trInvertibleSheaf。证明
\[
(\pi^*\mathscr{L}_1\cdots\pi^*\mathscr{L}_n)=\deg(X_1/X_2)(\mathscr{L}_1\cdots\mathscr{L}_n).
\]
（第一个\trIntersectionProduct{}在 $X_1$ 上，第二个在 $X_2$ 上。）提示：令 $d=\deg\pi$，并设 $d>0$。（$\pi$ 不\trDominant{}的情形另行处理，此时按约定 $d=0$，用\mbox{\trChevalley{}定理~8.4.2}。）论证：由\trIntersectionProduct{}的多重线性，只须处理诸 $\mathscr{L}_i$ 都是\trVeryAmple{}的情形。然后选取每个 $\mathscr{L}_i$ 的诸\trSection，使这些\trSection{}的一切\trIntersection{}都落在一个 $\pi$ 具有「真实次数 $d$」的\trOpenSubset{} $U$ 内。为找到 $U$：先用\mbox{练习~10.3.G} 找出一个稠密\trOpenSubset{} $\mathrm{U}'\subset X_2$，$\pi$ 在其上\trFinite。再用\mbox{有用练习~14.3.E} 证明：存在稠密\trOpenSubset{} $\mathrm{U}\subset\mathrm{U}'$，使得 $\pi_*\OO$ 在其上是一个\trRank{}为 $d$ 的\trLocallyFreeSheaf。用第 24 章的「平坦」语言来说，你是在证明：存在 $X_2$ 的一个稠密\trOpenSubset{} $\mathrm{U}$，$\pi$ 在其上既\trFinite{}又\trFlat（从而具有「恒常次数」；见\mbox{注记~24.3.8}）。（照例，把「射影」换成「固有」后结果也成立；见 [Ko1, Prop.~VI.2.11]。）

\textbf{20.1.8.}\ \emph{\trhRemark：更一般的\trProjection{}公式。}\quad 设 $\pi\colon X_1\to X_2$ 是诸\trProper{}\trVariety{}之间的一个\trProper\trMorphism，而 $\mathscr{F}$ 是 $X_1$ 上满足 $\dim\operatorname{Supp}\mathscr{F}\le n$ 的一个\trCoherentSheaf（于是由\mbox{练习~12.2.C}，$\dim\operatorname{Supp}\pi_*\mathscr{F}\le n$）。再设 $\mathscr{L}_1,\dots,\mathscr{L}_n$ 是 $X_2$ 上的\trInvertibleSheaf。那么
\[
(\pi^*\mathscr{L}_1\cdots\pi^*\mathscr{L}_n\cdot\mathscr{F})=(\mathscr{L}_1\cdots\mathscr{L}_n\cdot\pi_*\mathscr{F}).
\]
这个也称作\textbf{\trProjection{}公式}（并且以一种并不显然的方式推广了\mbox{练习~20.1.J}）。[FGIKNV, B.15] 与 [Ko1, Prop.~VI.2.11] 中给出了证明。由于我们不会用到这个版本的\trProjection{}公式，此处就不给出证明了。

\textbf{20.1.K.（练习：与丰富\trLineBundle{}\trIntersection）}\quad 设 $X$ 是一个射影 $k$-\trVariety，而 $\mathscr{L}$ 是 $X$ 上一个\trAmple{}\trLineBundle。证明：对 $X$ 的任意一个维数为 $n$ 的\trSubvariety{} $Y$，都有 $(\mathscr{L}^n\cdot Y)>0$。（提示：用\mbox{命题~20.1.3} 与\mbox{定理~16.2.2} 化归到 $\mathscr{L}$ 是\trVeryAmple{}的情形。然后证明：在由\trLinearSystem{} $|\mathscr{L}|$ 诱导的到\trProjectiveSpace{}的嵌入中，$(\mathscr{L}^n\cdot Y)=\deg Y$。）

\trNakai--\trMoishezon{}判别法（\mbox{定理~20.4.1}）断言：这一性质刻画了\trAmpleness。

\textbf{20.1.9.}\ \extrahard\ \emph{复射影情形中的\trCohomology{}诠释：推广\mbox{练习~18.4.H} 与 \S18.6.5。}\quad 若 $k=\mathbb{C}$，我们可以把 $(\mathscr{L}_1\cdots\mathscr{L}_n\cdot Y)$ 诠释为
\begin{equation*}\tag{20.1.9.1}\label{eq:chern-cap-class}
c_1((\mathscr{L}_1)_{\mathrm{an}})\cup\cdots\cup c_1((\mathscr{L}_n)_{\mathrm{an}})\cap[Y_{\mathrm{an}}]
\end{equation*}
在 $\mathrm{H}_0(Y_{\mathrm{an}},\mathbb{Z})$ 中的\trDegree。（回忆 $c_1((\mathscr{L}_i)_{\mathrm{an}})\in\mathrm{H}^2(X_{\mathrm{an}},\mathbb{Z})$，见\mbox{练习~18.4.H} 处的讨论。）证明它的一条途径，是利用\trIntersectionProduct{}与 \eqref{eq:chern-cap-class} 二者的多重线性，化归到 $\mathscr{L}_n$ 是\trVeryAmple{}的情形，于是 $\mathscr{L}_n\cong\OO(D)$，其中 $D$ 在 $Y$ 上的限制是一个有效\trCartierDivisor{} $E$。然后证明：若 $\mathscr{L}$ 是 $Y_{\mathrm{an}}$ 上的一个解析\trLineBundle{}（analytic line bundle），带有一个非零\trSection{} $E_{\mathrm{an}}$，则 $c_1(\mathscr{L})\cap[Y_{\mathrm{an}}]=[E_{\mathrm{an}}]$。最后，用对 $n$ 的归纳与\mbox{练习~20.1.D}。
